% MATH Companion Part VI / VI-31 / explanatory bridge only.
\paragraph{Codimension counts constraints along an irreducible subset.}
Let \(X\) be a Noetherian integral scheme, and let
\(Z\subseteq X\) be an irreducible closed subset with generic point
\(\eta_Z\).  The codimension of \(Z\) in \(X\) counts the
maximum length of chains of irreducible closed subsets joining
\(Z\) to \(X\), or equivalently the height of the prime
corresponding to \(\eta_Z\) on any affine neighbourhood of that
point:
\[
 \operatorname{codim}_X(Z)
 =\dim\mathcal O_{X,\eta_Z}
 =\operatorname{ht}(\mathfrak p_{\eta_Z}).
\]
This is an equality of \emph{local} invariants.  It does not
require us to assign a global coordinate equation to every closed
subset.

\paragraph{Example: one equation in the affine plane.}
Let \(X=\mathbb A^2_k\).  The irreducible line
\(L=V(y)\) is defined by the prime ideal \((y)\).
The chain \((0)\subsetneq(y)\) has length one; no nonzero
prime can occur strictly between them because \((y)\) has
height one.  Hence
\[
 \operatorname{codim}_{\mathbb A^2_k}(L)=1,
 \qquad \dim L=1,
 \qquad \dim\mathbb A^2_k=2.
\]
For the closed point \(P=V(x,y)\), the chain
\((0)\subsetneq(y)\subsetneq(x,y)\) has length two,
so \(\operatorname{codim}_X(P)=2\).  The two equations
\(x=0\) and \(y=0\) happen to be independent in this example,
but the height of an ideal cannot in general be inferred by just
counting its displayed generators.

\paragraph{Example: a coordinate line in three-space.}
Inside \(\mathbb A^3_k=\operatorname{Spec}k[x,y,z]\), let
\(Z=V(x,y)\), the \(z\)-axis.  Since
\[
 k[Z]=k[x,y,z]/(x,y)\cong k[z],
\]
we obtain \(\dim Z=1\).  The prime chain
\[
 (0)\subsetneq(x)\subsetneq(x,y)
\]
shows that \(\operatorname{ht}(x,y)\ge2\); the polynomial
ring dimension formula shows equality.  Therefore
\[
 \operatorname{codim}_{\mathbb A^3_k}Z=2,
 \qquad \dim Z+\operatorname{codim}_X Z=1+2=3.
\]
This calculation provides an algebraic meaning for the intuitive
phrase ``two independent conditions in three dimensions.''

\paragraph{A reducible closed set has several codimensions locally.}
Take \(Y=V(xy)\subseteq\mathbb A^2_k\).  It is the union
of the two coordinate axes.  Its irreducible components
\(V(x)\) and \(V(y)\) both have codimension one.  By the
standard convention for a nonempty closed set,
\[
 \operatorname{codim}_X Y=
 \min_{Z\text{ an irreducible component of }Y}
 \operatorname{codim}_X Z=1.
\]
The origin belongs to \(Y\), but its individual codimension in
\(X\) is two.  Thus \(\operatorname{codim}(Y,X)\) must not be
confused with the codimension of \emph{every} point in \(Y\).

\paragraph{When dimension plus codimension is justified.}
For an irreducible closed subset \(Z\) of an integral variety of
finite type over \(k\), the dimension formula gives
\[
 \dim Z+\operatorname{codim}_X(Z)=\dim X.
\]
For a nonempty closed \(Y\), taking the component of largest
dimension gives the analogous identity with the minimum-component
codimension defined above.  The formula can fail on arbitrary
non-equidimensional schemes.  When moving from varieties to
general schemes, one must check catenarity and the relevant
pure-dimensional hypotheses rather than quoting the identity
without explanation.
